\chapter{Анализ рынка и предметной области музыкального сопровождения заведений и мероприятий}
\label{ch:market}

\section{Музыкальное сопровождение заведений и мероприятий как отдельная сфера услуг}
\label{sec:market:music-service}

\subsection{Музыка как элемент клиентского опыта}

Музыкальное сопровождение коммерческих и событийных пространств представляет собой самостоятельную сферу услуг, находящуюся на пересечении индустрии развлечений, маркетинга и управления клиентским опытом. В заведениях общественного питания, гостиницах, барах, торговых и досуговых пространствах музыка выступает не только как эстетическое дополнение, но и как элемент целенаправленного проектирования атмосферы. Она задает эмоциональный фон посещения, влияет на восприятие уровня заведения, помогает формировать целостный образ пространства и поддерживает желаемую модель поведения клиентов. В этом смысле музыкальное оформление следует рассматривать как часть более широкой системы управления средой обслуживания, а не как случайный набор включенных композиций.

При этом важно разграничивать два близких, но функционально различных явления. В первом случае речь идет о фоновом музыкальном оформлении заведений. Такая музыка обычно не является самостоятельным предметом потребления: посетитель приходит в ресторан, кафе, бар, салон или гостиницу за основной услугой, а музыка сопровождает этот опыт, делая его более комфортным, цельным и соответствующим позиционированию места. Во втором случае музыка становится ядром самого события. Это характерно для концертов, клубных вечеринок, барных DJ-сетов и иных мероприятий, где музыкальная программа выступает одной из главных причин посещения, определяет сценарий вечера и организует коллективное внимание аудитории. Следовательно, фоновая музыка и музыкальное сопровождение событий различаются не только по уровню громкости или роли исполнителя, но прежде всего по месту музыки в ценностном предложении: в заведении она чаще поддерживает основной сервис, а на мероприятии во многом и является самим сервисом.

\subsection{Эволюция подходов: от фоновой музыки к дизайну среды}

Исторически представление о музыке как о специально проектируемом элементе коммерческой среды формировалось постепенно. Одним из наиболее известных ранних примеров стала компания Muzak, развивавшая централизованную поставку фоновой музыки для общественных, офисных и деловых пространств. Со временем название Muzak стало использоваться как нарицательное обозначение функциональной фоновой музыки, прежде всего ассоциируемой с «музыкой для лифтов». Однако в дальнейшем сама компания и рынок, который она олицетворяла, ушли от модели нейтрального звукового заполнения пространства к более сложному подходу, связанному с подбором музыкальных программ под тип бизнеса, аудиторию и желаемый эмоциональный эффект. В описании Д. Оуэна этот переход обозначается как движение к «audio architecture», то есть к проектированию аудиосреды как части общей архитектуры клиентского опыта~\cite{owen2006soundtrack}. Для настоящей работы этот исторический контекст важен по двум причинам. Во-первых, он показывает, что коммерческий спрос на управляемое музыкальное оформление существует давно. Во-вторых, он демонстрирует постепенный сдвиг от универсального фонового звучания к контекстному и сегментированному музыкальному подбору, что непосредственно связано с логикой современных рекомендательных систем.

\subsection{Эффекты музыкального оформления в индустрии сервиса}

Современные исследования описывают музыку в сервисных пространствах как элемент так называемого \textit{musicscape} — музыкального измерения среды обслуживания, воздействующего на посетителя через когнитивные и эмоциональные механизмы~\cite{trompeta2022music}. В рамках этой логики музыка может влиять на восприятие атмосферы, эмоциональное состояние, оценку качества пространства, удовлетворенность визитом и последующие поведенческие намерения. На уровне поведения в отдельных исследованиях анализируются такие показатели, как длительность пребывания, скорость потребления, расходы и готовность к повторному посещению. Вместе с тем академическая литература не дает оснований утверждать, что любое музыкальное сопровождение автоматически повышает выручку. Напротив, эффект зависит от типа пространства, контекста потребления, характеристик аудитории, акустических параметров и того, насколько музыка соответствует ситуации.

Эта осторожность особенно хорошо видна в мета-аналитических работах. В исследовании H. Roschk, S. M. C. Loureiro и J. Breitsohl, обобщившем результаты экспериментальной литературы по атмосферным стимулам, наличие музыки в среднем было связано с более высокими оценками удовольствия, удовлетворенности и поведенческих намерений потребителей~\cite{roschk2017atmospheric}. Однако более специализированный мета-анализ M.-A. Trompeta и соавторов, посвященный tourism and hospitality settings, показывает, что решающее значение имеет не сам факт присутствия музыки, а способ ее проектирования~\cite{trompeta2022music}. Авторы сопоставили 56 исследований и 209 эффектов и отдельно проанализировали пять измерений: наличие музыки, темп, громкость, конгруэнтность музыкального оформления контексту и субъективную привлекательность музыки для слушателя. Общий эффект простого наличия музыки оказался небольшим и лишь погранично значимым, тогда как предпочтительные и аффективные характеристики — \textit{music liking} и \textit{music congruence} — демонстрировали более выраженную связь с реакциями клиентов~\cite{trompeta2022music}.

Особенно значим для настоящей работы вывод о роли \textit{music liking}, то есть того, насколько звучащая музыка нравится посетителю. В мета-анализе Trompeta et al. этот фактор был связан с более позитивными когнитивными и эмоциональными реакциями, удовлетворенностью, намерением вернуться, готовностью рекомендовать заведение и рядом поведенческих показателей. В частности, для удовлетворенности был получен крупный эффект, а положительная связь также обнаружена для word-of-mouth intentions и repatronage intentions~\cite{trompeta2022music}. Сходный смысл имеет и \textit{music congruence} — соответствие музыки атмосфере места, категории услуги и ожидаемому поведению клиента. Конгруэнтная музыка в рассмотренных исследованиях ассоциировалась с более благоприятными оценками сервисной среды, когнитивными реакциями и поведенческими намерениями~\cite{trompeta2022music}.

Эти выводы подтверждаются и на уровне отдельных эмпирических исследований. В полевом эксперименте C. Caldwell и S. Hibbert в ресторанной среде музыкальное предпочтение посетителей оказалось более содержательным объясняющим фактором, чем темп композиций: музыка, которая нравилась гостям, была связана с более длительным пребыванием, более позитивной оценкой опыта, намерением вернуться и готовностью рекомендовать заведение~\cite{caldwell2002tempo}. При этом авторы подчеркивают, что связь между музыкальным оформлением и расходами не следует упрощать: важную роль играют промежуточные переменные, прежде всего длительность визита и эмоциональная оценка среды. Исследование N. Demoulin дополняет этот результат применительно к конгруэнтности музыкального оформления. Автор показывает, что соответствие музыки контексту сервисной среды влияет на эмоциональные реакции и когнитивные оценки, которые, в свою очередь, связаны с намерением повторного посещения~\cite{demoulin2011music}. Таким образом, для бизнеса более значимым оказывается не просто выбор условно «спокойной» или «быстрой» музыки, а способность обеспечить соответствие звучания аудитории и ситуации потребления.

Данный тезис имеет принципиальное значение для концепции, развиваемой в настоящей работе. Если позитивные эффекты музыкального сопровождения в значительной степени зависят от того, нравится ли музыка посетителям и воспринимается ли она как уместная, то персонализированный или адаптивный подбор музыкального потока может рассматриваться как механизм повышения качества аудиосреды. В случае стационарных плейлистов музыкальная программа формируется заранее и, как правило, ориентируется на усредненное представление о целевой аудитории заведения. Напротив, рекомендательная система, учитывающая предпочтения реально присутствующих посетителей, потенциально способна повысить вероятность совпадения музыкального потока с текущим составом аудитории. Речь идет не о гарантированном увеличении финансовых показателей, а о теоретически обоснованной прикладной гипотезе: более точное соответствие музыкального оформления вкусовому профилю гостей может усилить те положительные эффекты, которые в литературе уже связываются с \textit{liking} и \textit{congruence}.

Индустриальный интерес к подобному управлению музыкальной средой проявляется и в современной профессиональной дискуссии. Так, рассуждая о возможностях нейромузыки, заместитель руководителя сервиса «Яндекс Музыка» Е. Филиппов отметил, что ее в перспективе можно настраивать так, «чтобы люди покупали больше нужных товаров, сбрасывали больше веса в твоем спортзале, заказывали больше еды»~\cite{yandex_neuromusic_2022}. Эта реплика не является научным доказательством и не может использоваться как подтверждение причинно-следственного эффекта. Тем не менее она показательна как индустриальный комментарий: рынок уже рассматривает музыку не только как эстетический фон, но и как управляемый компонент пользовательского и коммерческого опыта.

\subsection{Музыкальное сопровождение мероприятий и функции диджея}

Отдельного рассмотрения требует музыкальное сопровождение мероприятий и вечеринок. В клубных форматах, на барных DJ-сетах и танцевальных событиях музыка выполняет не периферийную, а центральную сценарную функцию. Она определяет темп развития вечера, регулирует интенсивность коллективного переживания и формирует условия для вовлечения аудитории. В традиционной структуре клубного лайнапа можно выделить открывающего DJ, хедлайнера и закрывающего DJ. Задача opening DJ состоит, как правило, не в достижении пикового уровня энергии, а в постепенном «разогреве» аудитории, создании доверия к музыкальной атмосфере и подготовке зала к основному выступлению. Хедлайнер обеспечивает кульминационный участок события, тогда как closing DJ поддерживает или мягко снижает интенсивность опыта в финальной части вечера. Даже в барных форматах без крупного сценического события DJ нередко выполняет не только функцию воспроизведения музыки, но и роль живого куратора, реагирующего на состояние пространства.

Научные работы о восприятии ритма и \textit{groove} позволяют осторожно объяснить, почему музыка на мероприятиях способна усиливать вовлеченность. J. Stupacher, M. J. Hove и P. Vuust определяют \textit{groove} как приятное побуждение двигаться в такт музыкальному ритму и показывают, что оно связано с моторной активизацией, предсказуемостью ритмической структуры и балансом между повторяемостью и новизной~\cite{stupacher2023groove}. Авторы также подчеркивают социальный аспект совместного движения под музыку: синхронность действий в группе может усиливать чувство связанности между участниками~\cite{stupacher2023groove}. В отдельном исследовании Stupacher et al. показано, что культурная знакомость музыки и индивидуальный музыкальный вкус по-разному влияют на социальное сближение в условиях движения под музыку; при этом нравящаяся музыка особенно усиливала чувство близости при синхронном движении~\cite{stupacher2020familiarity}. Для контекста вечеринок это важно как косвенное обоснование гипотезы о том, что узнаваемый, нравящийся и ритмически вовлекающий материал может облегчать переход части аудитории от пассивного наблюдения к активному участию.

При этом клубный опыт нельзя свести к выбору «правильных» треков. Живой DJ взаимодействует с залом не только через музыкальный репертуар, но и через тайминг, драматургию сета, сценическое поведение и считывание реакции аудитории. Полевое исследование K. Rudnicki, T. Murrath и K. Poels в электронных танцевальных клубах показывает, что физическая вовлеченность DJ и отдельные аспекты его сценического поведения связаны с положительным аффектом аудитории~\cite{rudnicki2025djs}. Этот результат важен для интерпретации возможностей алгоритмических решений: рекомендательная система может поддерживать музыкальный выбор, но не исчерпывает всех функций живого исполнителя и не должна рассматриваться как его универсальная замена.

С учетом сказанного применение методов машинного обучения к музыкальному сопровождению заведений и мероприятий может быть обосновано в нескольких форматах. В повседневной деятельности кафе, ресторанов, баров и других B2B-клиентов рекомендательная система может выступать как виртуальный музыкальный куратор, формирующий поток, более релевантный текущей аудитории, чем статичные тематические плейлисты. В событийном контексте такая система может рассматриваться как потенциальный «виртуальный opening DJ», задача которого состоит в подборе начального музыкального материала для мягкого вовлечения аудитории и формирования комфортного уровня энергии. Наконец, система может использоваться как ассистент живого DJ: не заменяя профессиональное решение и чувство зала, она способна предоставлять дополнительную информацию о предполагаемых предпочтениях присутствующих гостей и тем самым помогать точнее выбирать композиции для разогрева аудитории.

Таким образом, музыкальное сопровождение заведений и мероприятий представляет собой самостоятельную сферу услуг, значение которой определяется не только распространенностью фоновой музыки в коммерческих пространствах, но и ее способностью участвовать в формировании клиентского опыта. Современная литература показывает, что влияние музыки на посетителей является контекстно зависимым, однако наиболее устойчивые положительные связи выявляются не для механического присутствия музыки и не только для ее физических параметров, а для субъективной привлекательности и уместности музыкального оформления. Именно этот вывод создает теоретическую основу для дальнейшего рассмотрения групповых музыкальных рекомендаций как инструмента повышения соответствия аудиосреды предпочтениям конкретной аудитории заведений и мероприятий.

\section{Сегментация рынка и существующие практики музыкального сопровождения заведений и мероприятий в России}
\label{sec:market:segmentation}

\subsection{Сегментация по типам заведений и функциям музыки}

Рынок музыкального сопровождения заведений и мероприятий в России неоднороден как по составу потенциальных клиентов, так и по тем функциям, которые музыка выполняет в конкретном пространстве. Наиболее очевидный способ его первичного описания — сегментация по типу бизнеса. Так, российский сервис «Звук Бизнес», работающий в сфере фоновой музыки для коммерческих пространств, выделяет среди целевых направлений HoReCa, ритейл, индустрию красоты, спорт, медицинские организации, общественные пространства и ряд смежных форматов. Внутри этих укрупнённых групп дополнительно обозначаются рестораны, кафе, бары, кофейни, гостиницы, магазины, торговые центры, салоны красоты, парикмахерские, барбершопы, фитнес-клубы, медицинские центры, аптеки, коворкинги и другие площадки~\cite{zvuk_business_site}.

Подобная отраслевая сегментация удобна для описания рынка фоновой музыки, поскольку в ней тип заведения действительно выступает важным ориентиром для выбора музыкального оформления. Кофейне, ресторану, фитнес-клубу и салону красоты, как правило, требуются различные темп, плотность звучания, эмоциональная окраска и режим обновления плейлистов. Однако для целей настоящего исследования одной отраслевой классификации недостаточно. Более содержательным является дополнительное разграничение по функции музыки в пространстве: в одних случаях она выступает атмосферным и сервисным элементом, а в других — становится центральной частью самого потребительского опыта.

В первом случае речь идёт о фоновой музыке для бизнеса. Её основная задача состоит не в том, чтобы быть самостоятельным предметом внимания посетителя, а в том, чтобы формировать нужную среду: поддерживать общее ощущение комфорта, соответствовать позиционированию заведения, усиливать впечатление от интерьера и не вступать в конфликт с основной услугой. Исследования атмосферных стимулов в коммерческих пространствах показывают, что музыка является значимым компонентом восприятия среды, однако эффект определяется не только фактом её присутствия, но и тем, насколько параметры звучания соответствуют контексту и ожиданиям аудитории~\cite{roschk2017atmospheric}. Аналогичный вывод сделан в мета-анализе исследований, посвящённых туризму и hospitality: авторы подчёркивают, что для клиентских реакций важен дизайн музыкального сопровождения, а предпочтительные характеристики музыки могут иметь более существенное значение, чем одни лишь физические параметры, например громкость или темп~\cite{trompeta2022music}.

Во втором случае музыка выполняет событийную функцию. Это характерно для вечеринок в барах, клубных форматов, концертных площадок, event-пространств и иных мероприятий, где музыкальное сопровождение не просто оттеняет происходящее, а участвует в создании самого продукта. В литературе по ночной экономике предлагается различать площадки, для которых музыкальные вечеринки являются основной заявленной активностью, и пространства, где музыка служит элементом дифференциации наряду с ресторанной или гостиничной функцией~\cite{ramon2022nightlife}. Такое различие особенно важно для данной работы: если в ресторане или салоне красоты музыка прежде всего дополняет основную услугу, то на вечеринке или в клубе она непосредственно влияет на структуру впечатления, вовлечённость гостей и динамику взаимодействия в пространстве.

Функции музыки, следовательно, существенно варьируются по сегментам. В ресторанах, кафе, кофейнях и барах в режиме обычного обслуживания она помогает задать характер атмосферы и поддержать соответствие между концепцией заведения и воспринимаемым образом пространства. Для части баров музыка может оставаться фоновым компонентом в дневное и ранневечернее время, но превращаться в ведущий элемент опыта в рамках специальных событийных программ. В индустрии красоты — салонах, парикмахерских, барбершопах и спа-пространствах — на первый план выходит комфортное, ненавязчивое и устойчивое по стилю музыкальное окружение, которое не отвлекает от процедуры и не создаёт раздражающего фонового эффекта. В гостиницах и аналогичных объектах важны длительный режим эксплуатации, согласованность музыкальной среды в разных зонах и возможность сохранять единый стандарт звучания на протяжении дня. Для фитнес-центров и спортивных студий музыка чаще связывается с поддержанием ритма активности и эмоциональной энергии занятий, что делает этот сегмент функционально особым по сравнению с пространствами спокойного обслуживания. Наконец, в клубах, на барных вечеринках, концертных и event-площадках музыка способна выступать ядром потребительского опыта и одним из ключевых факторов, определяющих восприятие всего события.

\subsection{Использование потребительских стриминговых сервисов и локальных музыкальных библиотек}

Такое функциональное различие принципиально и для анализа существующих рыночных практик. На практике бизнес решает задачу музыкального сопровождения несколькими способами, которые существенно различаются по стоимости, уровню профессионализации, степени управляемости и правовой модели. Наиболее простой путь — использование потребительских стриминговых сервисов, заранее собранных локальных плейлистов или подборок, которые формируются конкретным сотрудником, администратором или владельцем бизнеса. Такой подход обладает низким порогом входа: он не требует сложной инфраструктуры, быстро запускается и позволяет опираться на уже привычные музыкальные продукты. Однако он редко обеспечивает системное управление музыкальной средой, слабо масштабируется на сети заведений и, что особенно важно для логики настоящей работы, никак не связан с фактическим составом аудитории, находящейся в помещении в конкретный момент времени.

Кроме того, использование потребительских музыкальных сервисов в публичной коммерческой среде связано с правовыми рисками. Например, в условиях использования сервиса «Яндекс Музыка» прямо указано, что сервис и материалы предназначены для личного некоммерческого использования~\cite{yandex_music_terms}. В рамках данного подраздела этот вопрос не рассматривается подробно, поскольку юридические аспекты публичного воспроизведения музыки будут анализироваться отдельно далее. Тем не менее сама граница между потребительской и коммерческой моделью важна для понимания того, почему часть бизнеса переходит к специализированным B2B-решениям.

\subsection{B2B-сервисы и практики брендированного музыкального оформления}

Отдельную практику образует формирование брендированной музыкальной среды. В этом случае музыка используется не только как фон, но и как часть идентичности заведения или сети. Российским примером является проект Surf Coffee Propaganda Machine, в рамках которого бренд развивает собственное самостоятельное музыкально направление и публикует регулярные сеты от своих штатных диджеев~\cite{surf_soundcloud}. В более широком теоретическом контексте подобные практики соотносятся с подходами sonic branding, где звук рассматривается как выражение идентичности бренда и способ выстраивания узнаваемого аудиального образа в повседневном потребительском опыте~\cite{khamis2021sonic}.

Показательно, что музыкальная айдентика отдельных ритейл-брендов может получать самостоятельную жизнь за пределами торгового пространства. Так, архивирование и переосмысление внутримагазинных плейлистов GAP стало заметным онлайн-феноменом: вокруг этих подборок сложилось сообщество заинтересованных пользователей, а сам бренд впоследствии использовал эту культурную память в коммуникационных активностях~\cite{vogue_gap_playlists_2022}. Похожим образом на стриминговых платформах возникают пользовательские плейлисты, стилизованные под музыку, звучащую в магазинах ZARA; некоторые из них собирают тысячи сохранений, что указывает на самостоятельный интерес аудитории к узнаваемой музыкальной атмосфере бренда~\cite{spotify_zara_store_playlist_2026}.

Для темы нашей работы данный аспект двойственное значение. С одной стороны, он показывает, что музыкальное оформление может быть не вспомогательным, а ценным элементом брендового опыта, способным удерживаться в памяти посетителей и воспроизводиться ими самостоятельно. С другой стороны, при столь выраженной музыкальной айдентике прямая персонализация «под каждого текущего посетителя» может вступать в противоречие с целостностью брендового образа. Поэтому применительно к ритейлу и брендам с ярко выраженным музыкальным стилем более перспективной представляется не неограниченная персонализация, а рекомендательная модель, работающая внутри заранее заданных жанровых, эмоциональных и стилистических рамок.

Ещё одна распространённая рыночная модель — специализированные B2B-сервисы музыки для бизнеса. Их ценностное предложение обычно строится вокруг нескольких характеристик: легальная модель использования контента, наличие каталогов и тематических подборок, инструменты централизованного или удалённого управления воспроизведением, а также готовые музыкальные решения, адаптированные к различным типам заведений~\cite{zvuk_business_site}. Для бизнеса это позволяет перейти от спонтанного и персонально зависимого выбора музыки к более устойчивой операционной модели. Однако сама логика таких решений чаще ориентируется на формат бизнеса, категорию целевой аудитории или заранее настроенный профиль пространства, а не на динамически меняющийся состав реальных посетителей в конкретной точке. Следовательно, даже профессионализированная B2B-модель не устраняет полностью исследовательский разрыв, на который направлена настоящая работа.

\subsection{Музыкальное оформление мероприятий}

В событийном сегменте исторически сохраняется другая практика — ручной музыкальный отбор, работа DJ, участие музыкальных редакторов, технических специалистов и использование профессионального программного обеспечения для подготовки и исполнения сетов. Современные DJ-платформы прямо позиционируются как инструменты для организации музыкальных библиотек, подготовки выступлений, быстрой адаптации треков и управления логикой сетов в ходе живого исполнения~\cite{rekordbox_software}. Такая модель особенно характерна для клубов, барных вечеринок, концертных мероприятий и event-площадок, где музыкальное решение должно учитывать развитие ситуации «здесь и сейчас». В отличие от фоновой музыки для бизнеса, событийное музыкальное сопровождение уже в большей степени направлено на наблюдение за реакцией аудитории и оперативную коррекцию репертуара. Однако эта адаптация, как правило, реализуется через профессиональную интуицию, визуальную оценку публики и опыт конкретного исполнителя, а не через автоматизированный учёт структуры присутствующей аудитории и её музыкальных предпочтений.

Таким образом, современный рынок музыкального сопровождения заведений и мероприятий можно описать как совокупность нескольких сосуществующих практик. На одном полюсе находятся локальные плейлисты и потребительские музыкальные продукты, обеспечивающие минимальный порог входа, но дающие низкую управляемость и ограниченную воспроизводимость результата. Далее располагаются брендированные музыкальные решения, в которых звучание становится частью фирменного опыта и инструментом коммуникации. Следующую группу составляют специализированные B2B-сервисы, предлагающие более формализованную инфраструктуру управления фоновой музыкой. Наконец, в событийном сегменте доминируют профессиональные формы ручной адаптации — DJ-сопровождение, музыкальная редактура и работа с программными инструментами живого выступления.

В большинстве случаев музыкальное оформление ориентируется либо на тип бизнеса, либо на заранее сформулированный брендовый стиль, либо на экспертный вкус конкретного куратора или DJ. При этом динамический состав фактической аудитории пространства обычно не выступает самостоятельным параметром управления музыкальной средой. Именно эта сфера создаёт основание для постановки исследовательской задачи нашей работы: проверить, насколько перспективна модель коллективных музыкальных рекомендаций, способная учитывать предпочтения реальных посетителей заведения или мероприятия при сохранении заданных контекстных ограничений.

\subsection{Приоритетные сегменты для дальнейшего исследования}

Исходя из этой логики, в рамках дальнейшего исследования целесообразно сфокусироваться не на всём рынке музыкального сопровождения, а на тех сегментах, где связь между составом аудитории и качеством музыкального опыта потенциально наиболее значима. При этом ряд направлений не рассматривается как приоритетный фокус. К ним относятся фитнес-центры, поскольку их клиентская экономика абонементов, не всегда поощряет повторные посещения и черезмерное удержание пользователей в залах. Также вне основного фокуса остаётся значительная часть ритейла: хотя музыкальная айдентика магазина может быть важной частью восприятия бренда, персонализация под текущий поток посетителей в этом случае не всегда очевидно совместима с задачей сохранения стабильного фирменного образа.

Основной исследовательский интерес представляют четыре группы сегментов. Во-первых, это мероприятия и пространства, где музыка является ядром пользовательского опыта: барные вечеринки, клубы, концертные и event-площадки, а также организаторы событий. Для таких форматов ценность музыкального решения максимально связана с коллективной реакцией присутствующих и с тем, насколько репертуар способствует вовлечённости аудитории.

Во-вторых, это заведения с фоновой музыкой, где она заметно влияет на восприятие пространства и клиентский опыт: бары вне событийного режима, рестораны, кафе и кофейни. Здесь особенно важно исследовать возможность мягкой адаптации музыкальной среды без разрушения заданного образа заведения.

В-третьих, это сфера услуг, в которой клиент проводит в пространстве заметное время: салоны красоты, парикмахерские и барбершопы. Данный выбор согласуется с результатами исследований атмосферных эффектов музыки: мета-анализ Roschk et al. указывает на тенденцию к более сильным эффектам музыкальных стимулов в сервисных пространствах по сравнению с ритейлом~\cite{roschk2017atmospheric}, а Trompeta et al. показывают значимость качественно спроектированного музыкального сопровождения в hospitality-контекстах~\cite{trompeta2022music}.

В-четвёртых, перспективным представляется сегмент развлечений — боулинги, бильярдные, развлекательные центры и иные пространства досуга, где знакомая и эмоционально подходящая музыка может быть важной частью общей вовлечённости, хотя степень и характер этого эффекта требуют отдельной проверки в рамках прикладной модели.

Таким образом, выбранная сегментация позволяет перейти от широкого описания рынка музыкального сопровождения к более точной постановке предмета исследования. Для ВКР принципиально важны не любые форматы использования музыки в бизнесе, а те пространства, где музыкальная среда влияет на восприятие опыта, имеет потенциал для адаптации и при этом сегодня, как правило, не учитывает динамический состав фактической аудитории. Именно в этих сегментах концепция коллективных музыкальных рекомендаций может быть теоретически и практически наиболее содержательной.

\section{Конкурентный анализ существующих решений на рынке музыки для бизнеса}
\label{sec:market:competitors}

Российский рынок музыки для бизнеса в открытых источниках выглядит не как единый прозрачный рынок с общепринятыми долями, а как совокупность нескольких моделей. С точки зрения постановки настоящей ВКР ключевое конкурентное различие проходит не между «больше треков» и «меньше треков», а между тремя ценностями продукта. Первая~--- юридическая безопасность и снижение административной нагрузки. Вторая~--- аудиомаркетинг и музыкальный брендинг. Третья~--- вовлечение гостя в формирование эфира. Поэтому в рамках конкурентного анализа целесообразно отдельно рассмотреть классические сервисы музыки для бизнеса и сервисы заказа треков гостями.

\subsection{Сервисы музыки для бизнеса}

Первый класс решений составляют сервисы с собственным или прямолицензированным каталогом, которые позиционируют подписку как достаточную для законного публичного воспроизведения без дополнительных выплат в РАО и ВОИС. К этой группе относятся, например, Cubic Music, Бубука, Mubicloud, IMSTREAM, Aimon, RadioSparx, а также часть предложений МузКафе и Маркет Радио~\cite{cubic_music_site,bubuka_site,mubicloud_site,imstream_site,aimon_site,radiosparx_site,muscafe_site,market_radio_site}. Второй близкий класс образуют сервисы с более широким каталогом и сопутствующей работой с отчетностью и правами, где показательным примером является FONMIX, а также премиальные конфигурации Звук Бизнес~\cite{fonmix_site,zvuk_business_site}. В обоих случаях основная ценность для клиента состоит в легальном музыкальном фоне, централизованном управлении воспроизведением, аудиомаркетинге, брендированных плейлистах и снижении операционной нагрузки на персонал.

Отдельно следует выделить уже существующий продукт «Яндекс Станция. Музыка для бизнеса». Он реализован как совместный навык Яндекса и Cubic Music, доступный через Яндекс Станцию или ТВ Станцию, и предлагает бизнесу набор готовых плейлистов, голосовое управление, личный кабинет объектов и сценарии умного дома~\cite{yandex_station_business}. Для настоящей работы этот продукт важен как подтверждение того, что у Яндекса уже есть B2B-направление в музыкальном сопровождении бизнеса. При этом по публичному описанию он остается ближе к модели готовых плейлистов и не решает задачу динамической адаптации музыкального потока под фактический состав посетителей.

\begingroup
\small
\singlespacing
\setlength{\tabcolsep}{3pt}
\begin{longtable}{p{0.15\textwidth}p{0.18\textwidth}p{0.20\textwidth}p{0.22\textwidth}p{0.17\textwidth}}
\caption{Сравнительная матрица B2B-сервисов музыки для бизнеса}
\label{tab:business-music-competitors}\\
\toprule
Сервис & Стоимость / модель & Основные функции & Правовая модель и отчетность & Отличие от проектируемого сервиса \\
\midrule
\endfirsthead
\toprule
Сервис & Стоимость / модель & Основные функции & Правовая модель и отчетность & Отличие от проектируемого сервиса \\
\midrule
\endhead
Cubic Music / Кубик Медиа & Цена основного интеграционного сервиса публично не фиксирована; совместный продукт с Яндексом оплачивается помесячно & Брендовые плейлисты, музыкальное расписание, управление филиалами, внутреннее радио, аудиоролики & Акцент на собственный каталог и легальное воспроизведение без дополнительных платежей; публичный экспорт отчетов в РАО/ВОИС не является ядром коммуникации & Сильный аудиомаркетинг, но без групповой адаптации под посетителей \\
\addlinespace
Яндекс Станция. Музыка для бизнеса & Публично указана месячная подписка за устройство при годовой оплате & 30+ плейлистов, голосовое управление, личный кабинет объектов, сценарии умного дома & Легальное воспроизведение через навык Cubic; отчетность в РАО/ВОИС публично не заявлена & Близкая экосистема, но модель готовых плейлистов, а не коллективных рекомендаций \\
\addlinespace
Звук Бизнес & Стандартная цена зависит от числа точек; есть тестовый период и отдельные промопредложения & Волны под тип бизнеса, web/mobile/player, аудиореклама, видеоэкраны, премиальное кураторство & Легальная библиотека; в премиальной модели возможны отдельные договоры с РАО/ВОИС и подготовка отчетов & Закрывает управление эфиром и аудиорекламу, но не делает участие гостей ядром продукта \\
\addlinespace
Бубука & Публичные тарифы начинаются с подписки для одной точки & Готовые и индивидуальные плейлисты, офлайн-режим, рекламный кабинет, QR-инструменты & Позиционируется как музыка без дополнительных выплат; есть детальная отчетность по трекам & Фокус на легальности и операционном управлении, а не на групповом вкусовом профиле аудитории \\
\addlinespace
FONMIX & Публичные тарифы по уровням; права могут оплачиваться отдельно & Музыка и видеооформление, конструктор плейлистов, аудио- и видеореклама, мониторинг & Один из наиболее развитых reporting-first подходов; сервис не всегда заменяет РАО/ВОИС & Сильная отчетность, но персонализация по фактическим посетителям не является основной функцией \\
\addlinespace
RadioSparx & Публичные тарифы зависят от типа подписки и каталога & Станции, плейлисты, автоматизация смены музыки, собственная станция & Direct-license модель; отчеты формируются для правообладателей & Решает задачу легального фона, но не использует историю слушателей конкретной аудитории \\
\addlinespace
Mubicloud & Публичные тарифы от базовой подписки; есть расширенные планы & Облачный плеер, готовые сборники, персонализированные плейлисты, аудиомаркетолог & Заявляется замена платежей в РАО/ВОИС за счет собственной правовой модели & Персонализация относится к бренду/точке, а не к группе текущих посетителей \\
\addlinespace
МузКафе & Публично указаны тарифы от базовой подписки и индивидуальные решения & Музыкальные серверы, аудиобрендинг, рекламные сообщения, сопровождение прав & Возможна модель собственной медиатеки и сопровождение клиента по документам & Интеграционный и правовой аутсорсинг без social-recommendation механики \\
\addlinespace
IMSTREAM & Публичные тарифы от базового уровня; пробный период & Лицензионная музыка, готовые подборки, возможность работы без интернета & Подчеркивается отсутствие дополнительных выплат и юридическая поддержка & Закрывает легальный фон, но не предлагает динамический учет аудитории \\
\addlinespace
Аймон & Публичные тарифы по уровням & Станции, расписание музыки и контента, рекламный менеджер, browser-based управление & Легальная музыка с документами; явный экспорт РАО/ВОИС не является центральным элементом & Простая SMB-модель без коллективной персонализации \\
\addlinespace
Маркет Радио & Бесплатные и платные уровни; стоимость зависит от набора услуг & Онлайн-радио, индивидуальные программы, удаленное управление, аудиоролики & Сочетаются прямая лицензия и калькуляция платежей; автоматическая отчетность публично не выделена & Ближе к радио и аудиомаркетингу, чем к рекомендательной системе \\
\bottomrule
\end{longtable}
\endgroup

Сравнительная матрица в Таблице~\ref{tab:business-music-competitors} показывает, что большинство игроков уже закрывает базовую потребность бизнеса в легальном и управляемом музыкальном фоне. Поэтому проектируемый сервис не должен позиционироваться как еще один «плеер для бизнеса» или как прямой аналог существующих B2B-подписок. Его принципиальное отличие состоит в источнике персонализации: музыка подбирается не только по типу заведения, расписанию или заранее заданному брендовому профилю, а с учетом предпочтений фактически присутствующей группы посетителей. В этом смысле классические сервисы музыки для бизнеса и проектируемый сервис частично обращены к одному B2B-клиенту, но решают разные задачи. Первые обеспечивают легальный каталог, воспроизведение, управление точками и аудиомаркетинг; проектируемый сервис добавляет рекомендательный слой, который потенциально может работать поверх такой инфраструктуры и отвечать за динамическую адаптацию потока к текущей аудитории.

\subsection{Сервисы заказа музыки посетителями}

Отдельную группу образуют сервисы заказа музыки гостями, где главным ценностным предложением становятся QR-механика, очередь треков, VIP- или bidding-логика и дополнительная выручка для заведения. К этой группе относятся GUSLI FM и «Бит за Бид»~\cite{gusli_site,bitzabid_site}. В отличие от классических B2B-сервисов фоновой музыки, такие продукты делают участие посетителя явной частью пользовательского сценария: гость открывает интерфейс, выбирает конкретную композицию и тем самым получает возможность повлиять на эфир.

\begingroup
\small
\singlespacing
\setlength{\tabcolsep}{3pt}
\begin{longtable}{p{0.13\textwidth}p{0.17\textwidth}p{0.22\textwidth}p{0.18\textwidth}p{0.22\textwidth}}
\caption{Сравнительная матрица сервисов заказа музыки гостями}
\label{tab:ordering-competitors}\\
\toprule
Сервис & Стоимость / модель & Основные функции & Правовая модель и отчетность & Отличие от проектируемого сервиса \\
\midrule
\endfirsthead
\toprule
Сервис & Стоимость / модель & Основные функции & Правовая модель и отчетность & Отличие от проектируемого сервиса \\
\midrule
\endhead
GUSLI FM & Для заведения может быть бесплатным; монетизация строится вокруг заказов гостей & QR/app-заказы, VIP-песня, блокировки, реклама, история заказов, кэшбэк & Требуются договоры с РАО и ВОИС; есть выгрузка отчетов & Единоличные заказы гостей, а не мягкое коллективное формирование потока \\
\addlinespace
Бит за Бид & Бесплатно для заведений; комиссия с заказов & QR-заказы, аукцион треков, интерфейс DJ, возврат средств при отклонении & Публичная коммуникация сосредоточена на платежной и event-механике; отчетность РАО/ВОИС не раскрыта & Сильная интерактивность, но влияние зависит от активности и платежеспособности отдельных гостей \\
\bottomrule
\end{longtable}
\endgroup

Сервисы, представленные в Таблице~\ref{tab:ordering-competitors}, ближе к механике visitor ordering, чем к рекомендательной системе в строгом смысле. Они конкурируют не столько за качество автоматического подбора музыкального потока, сколько за сценарий вовлечения гостя и за дополнительную монетизацию конкретного заказа. В таких моделях музыкальный выбор становится дискретным и индивидуализированным: отдельный посетитель выбирает или оплачивает трек, а заведение затем вынуждено управлять очередью, блокировками, модерацией и возможными конфликтами между запросами гостей и музыкальной политикой площадки.

Проектируемый сервис принципиально отличается от этой группы конкурентов. Он не предполагает продажи права поставить конкретную композицию и не строит влияние пользователя на платежной или аукционной логике. Вклад посетителя является не прямым заказом трека, а частью агрегированного группового профиля, сформированного на основе музыкальных предпочтений присутствующих пользователей и ограничений, заданных заведением. Поэтому проектируемый сервис не конкурирует напрямую с GUSLI FM и «Бит за Бид» как с инструментами монетизации одиночных заказов. Он решает другую задачу: не передать эфир наиболее активному или платежеспособному гостю, а мягко адаптировать общий музыкальный поток к фактическому составу аудитории без разрушения сценария заведения.

Главный рыночный разрыв можно сформулировать следующим образом: существующие сервисы закрывают легальность, каталог, управление музыкой, аудиомаркетинг и иногда интерактивный заказ треков, но почти не предлагают динамическую адаптацию музыкального потока под фактический состав посетителей. Именно этот разрыв формирует продуктовую нишу для сервиса коллективных музыкальных рекомендаций.

\section{Правовые аспекты публичного воспроизведения музыки и требования к отчетности}
\label{sec:market:legal}

\subsection{Правовая природа публичного воспроизведения и роль РАО/ВОИС}

Для легального публичного воспроизведения музыки необходимо наличие правового основания на соответствующее использование объектов интеллектуальной собственности. Согласно Гражданскому кодексу РФ, публичное исполнение произведения представляет собой его представление в живом исполнении или с помощью технических средств в месте, открытом для свободного посещения, либо в месте, где присутствует значительное число лиц, не принадлежащих к обычному кругу семьи~\cite{gk_rf_1270}. Следовательно, воспроизведение музыкальных композиций в кафе, ресторанах, магазинах, салонах, гостиницах и иных коммерческих пространствах подпадает под режим публичного исполнения и требует урегулирования прав с правообладателями.

Распространенным механизмом такого урегулирования являются организации по коллективному управлению правами, которые осуществляют сбор вознаграждения в пользу правообладателей за различные виды использования произведений, исполнений и фонограмм. В России в рассматриваемой сфере действуют две ключевые аккредитованные организации: Российское авторское общество, управляющее авторскими правами на музыкальные произведения при их публичном исполнении, и Всероссийская организация интеллектуальной собственности, осуществляющая сбор вознаграждения за публичное исполнение фонограмм, опубликованных в коммерческих целях~\cite{rao_usage,vois_about}.

Особенность государственной аккредитации состоит в том, что в предусмотренных законом сферах соответствующая организация вправе действовать не только в интересах правообладателей, заключивших с ней договор, но и в отношении правообладателей, которые не передали ей полномочия напрямую, если они специально не изъяли свои права из управления. Такой механизм закреплен в статье 1244 ГК РФ и позволяет пользователям получать разрешение или урегулировать выплату вознаграждения через одну аккредитованную организацию, а не выстраивать индивидуальные отношения с каждым автором, издателем, исполнителем или изготовителем фонограммы~\cite{gk_rf_1244}.

Применительно к воспроизведению музыкального трека в заведении необходимо учитывать двойственную правовую природу используемого контента. Во-первых, используется само музыкальное произведение — объект авторских прав, что предполагает получение разрешения на его публичное исполнение, как правило, через заключение договора с РАО~\cite{gk_rf_1270,rao_usage}. Во-вторых, используется фонограмма, то есть конкретная звукозапись произведения. В соответствии со статьей 1326 ГК РФ публичное исполнение фонограммы, опубликованной в коммерческих целях, допускается без отдельного разрешения изготовителя фонограммы и исполнителя, но при условии выплаты им вознаграждения; сбор такого вознаграждения осуществляется через аккредитованную организацию, которой в данной сфере выступает ВОИС~\cite{vois_about,gk_rf_1326}. Поэтому для правомерного использования массового музыкального репертуара в публичном пространстве бизнесу, как правило, требуется одновременно урегулировать отношения с РАО и ВОИС.

Как было показано в предыдущем подразделе, значительная часть сервисов музыки для бизнеса предлагает каталоги, которые позиционируются как пригодные для публичного воспроизведения без дополнительных отчислений в РАО и ВОИС. Однако такая модель обычно основана на использовании музыкального контента, права на который урегулированы напрямую самим сервисом. Подобные каталоги, как правило, отличаются от репертуара массовых стриминговых платформ и в меньшей степени включают популярную и широко узнаваемую музыку. На это косвенно указывает наличие у ряда B2B-сервисов отдельных опций, тарифов или решений, ориентированных именно на доступ к известному коммерческому репертуару.

Отдельно необходимо отметить, что использование кавер-версий популярных композиций само по себе не устраняет обязанности по урегулированию авторских прав на музыкальное произведение. Даже если в заведении воспроизводится не оригинальная запись, а новая фонограмма, созданная другим исполнителем, при этом продолжает использоваться исходное музыкальное произведение, права на которое сохраняются за его авторами или иными правообладателями~\cite{gk_rf_1270,gk_rf_1260}. Аналогичным образом ремиксы и иные переработки музыкального материала относятся к производным произведениям; при их использовании подлежат защите как права автора переработки, так и права авторов произведений, на которых такая переработка основана~\cite{gk_rf_1260}.

\subsection{Юридические риски бездоговорного использования музыки}

Наиболее существенным практическим риском для собственников бизнеса, использующих музыку без надлежащего оформления отношений с РАО и ВОИС, является предъявление гражданско-правовых требований о взыскании компенсации. В судебной практике факты публичного исполнения музыкальных произведений и фонограмм нередко фиксируются представителями правообладателей или организаций по коллективному управлению правами посредством посещения заведения и проведения аудио- или видеофиксации воспроизводимого контента. Впоследствии такие материалы используются при направлении претензии и в ходе судебного разбирательства~\cite{sudact_a19_2521_2025}.

Гражданский кодекс РФ предусматривает возможность взыскания компенсации за нарушение исключительных прав на произведение и на объекты смежных прав. По статье 1301 ГК РФ компенсация за нарушение прав на произведение может составлять от 10 000 рублей до 10 миллионов рублей, а по статье 1311 ГК РФ аналогичный диапазон установлен для нарушения исключительных прав на фонограмму и иные объекты смежных прав~\cite{gk_rf_1301,gk_rf_1311}. При этом размер требований может определяться применительно к каждому отдельному объекту нарушения. Поэтому воспроизведение одной музыкальной композиции потенциально затрагивает одновременно права на произведение и права на фонограмму, а совокупный размер требований увеличивается пропорционально количеству зафиксированных треков.

Так, в условной ситуации, когда истцом зафиксировано публичное воспроизведение десяти музыкальных треков, минимальный размер требований при одновременном предъявлении претензий по авторским и смежным правам может составить 200 000 рублей: по 10 000 рублей за нарушение прав на произведение и по 10 000 рублей за нарушение прав на фонограмму в отношении каждого трека. Данный пример демонстрирует не конкретный исход судебного дела, а принцип формирования потенциального размера компенсации на основе установленных законом минимальных пределов~\cite{gk_rf_1301,gk_rf_1311}.

Таким образом, заключение договоров с РАО и ВОИС является для бизнеса не только средством соблюдения требований законодательства, но и способом снижения предсказуемых юридических и финансовых рисков. Наличие таких соглашений переводит музыкальное сопровождение заведения в правомерную плоскость, создает формализованный порядок выплаты вознаграждения правообладателям и снижает вероятность предъявления требований, связанных с бездоговорным публичным воспроизведением музыкального контента.

\subsection{Требования к отчетности}

Статья 1243 ГК РФ устанавливает, что пользователи по требованию организации по управлению правами на коллективной основе обязаны представлять отчеты об использовании объектов авторских и смежных прав. Содержательно отчетность предполагает фиксацию информации о фактически использованных объектах. РАО публикует формы отчетов об использовании произведений, а в формах для публичного исполнения указываются сведения о названии произведения, исполнителе, авторе музыки, авторе текста и иных параметрах, необходимых для идентификации использованного репертуара~\cite{rao_report_forms}. Формы ВОИС для отчетов о публичном исполнении фонограмм предусматривают отражение наименования фонограммы, исполнителя, автора музыки, автора текста, изготовителя фонограммы и количества исполнений~\cite{vois_report_form}. Генерация подобных отчетов должно быть неотъемлемой частью сервиса музыки для бизнеса.

\subsection{Правовая неопределенность источника фонограмм}

Однако даже при наличии выстроенного механизма лицензирования публичного исполнения через РАО и ВОИС сохраняется отдельный правовой вопрос: из какого источника бизнес получает фонограммы для фактического воспроизведения. Договоры с организациями по коллективному управлению правами урегулируют выплату вознаграждения за публичное исполнение произведений и фонограмм, но ому каталогу и не заменяют соглашение с цифровой платформой, через которую осуществляется воспроизведение.

В этом отношении потребительские музыкальные стриминги не могут рассматриваться как полноценный инструмент легального музыкального сопровождения бизнеса. Так, действующие Условия использования сервиса «Яндекс Музыка» прямо устанавливают, что сервис, все размещенные в нем материалы, плеер и база данных предназначены исключительно для личного некоммерческого использования. Любое использование сервиса и музыкальных материалов в коммерческих целях без отдельного согласования с правообладателем сервиса запрещается~\cite{yandex_music_terms}. Следовательно, даже если заведение заключило договоры с РАО и ВОИС, воспроизведение музыки через обычную пользовательскую подписку на стриминговый сервис остается несоответствующим условиям самой цифровой платформы.

В результате возникает разрыв между правовой моделью публичного исполнения и реальными техническими сценариями, к которым прибегает бизнес. Если организация не использует специализированный B2B-сервис, наиболее доступными источниками фонограмм становятся либо личные аккаунты обычных стриминговых платформ, либо локальные аудиофайлы, происхождение и правовой статус которых могут быть неопределенными. Первый вариант позволяет технически воспроизводить популярный музыкальный репертуар, но противоречит условиям потребительской лицензии стриминга. Второй вариант не связан с нарушением пользовательского соглашения платформы, однако не решает вопрос о законности получения самих файлов и наличии у пользователя необходимых прав на их дальнейшее использование.

Эта ситуация имеет принципиальное значение для проектируемого сервиса групповых музыкальных рекомендаций. Если такой продукт предполагает работу с популярной музыкой, сопоставимой по составу с каталогом массового стриминга, то правовой проблемой становится не только публичное исполнение как таковое, но и наличие у платформы самостоятельных B2B-прав на соответствующий музыкальный репертуар. Иными словами, для создания юридически устойчивого сервиса недостаточно обеспечить заключение договоров клиентов с РАО и ВОИС; необходимо также получить от правообладателей, прежде всего музыкальных лейблов и иных владельцев прав на фонограммы, объем разрешений, позволяющий использовать соответствующий каталог в коммерческих пространствах. Без такого расширенного лицензионного контура сервис фактически оказывается зависимым от потребительской инфраструктуры стриминга, которая юридически предназначена для иных целей.

После сокращения доступности части зарубежного музыкального репертуара в российских стриминговых сервисах в 2022 году, когда отдельные иностранные правообладатели начали отзывать права на размещение контента в России~\cite{forbes_rights_withdrawal_2022}.

Показательным в данном контексте является механизм пользовательской загрузки фонограмм, предусмотренный в действующих Условиях использования «Яндекс Музыки». Сервис предоставляет авторизованному пользователю возможность размещать фонограммы как для собственного прослушивания, так и с возможностью прослушивания другими пользователями. При этом пользователь самостоятельно несет ответственность перед третьими лицами за размещаемые материалы, гарантирует отсутствие нарушения прав третьих лиц и обязуется самостоятельно урегулировать возможные претензии правообладателей; сам сервис, в свою очередь, прямо указывает, что не обязан предварительно проверять пользовательские материалы на соответствие законодательству об авторском праве и смежных правах~\cite{yandex_music_terms}. Данная конструкция демонстрирует модель, при которой цифровая платформа предоставляет техническую возможность работы с пользовательским контентом, но не принимает на себя полный объем правовых гарантий относительно происхождения каждой загружаемой фонограммы.

Таким образом, ключевая проблема состоит в том, что в России сосуществуют, но не совпадают между собой три правовых контура: договоры с РАО и ВОИС, регулирующие выплату вознаграждения за публичное исполнение; пользовательские соглашения массовых музыкальных стримингов, ограничивающие использование сервиса личными некоммерческими целями; и специальные B2B-лицензии на музыкальные каталоги, необходимые для законного коммерческого воспроизведения контента в заведениях.

\subsection{Возможные модели правового оформления сервиса}

Следовательно, при проектировании сервиса групповых музыкальных рекомендаций можно выделить несколько возможных моделей его правового и организационного оформления. При этом ключевым ограничением остается то, что сам сервис в базовой концепции может позиционироваться не как поставщик музыкального контента, а как информационная система, формирующая рекомендации для бизнеса на основе предпочтений посетителей.

Первый вариант предполагает использование сервиса в связке с режимом личного прослушивания треков из библиотеки «Яндекс Музыки». В этом случае система формирует рекомендованный плейлист, а воспроизведение фактически осуществляется через пользовательский аккаунт стримингового сервиса. Подобная практика в настоящее время широко распространена в заведениях, несмотря на то что пользовательские соглашения музыкальных стримингов, как правило, предусматривают использование контента исключительно в личных, а не коммерческих целях. Таким образом, данный вариант не создает полноценной юридически прозрачной модели публичного воспроизведения, но опирается на уже сложившуюся рыночную практику, при которой бизнесы используют потребительские музыкальные сервисы для фонового оформления пространства.

Второй вариант состоит в создании режима плеера, при котором проектируемый сервис предоставляет программное обеспечение для воспроизведения музыкальных файлов, загруженных самим пользователем. В такой модели ответственность за происхождение аудиофайлов и наличие необходимых прав на их использование переносится на бизнес-клиента. Сервис, в свою очередь, может сопоставлять загруженные композиции с рекомендованными треками при помощи аудиоотпечатков, формировать на этой основе персонализированные и групповые подборки, а также обеспечивать автоматизированное воспроизведение. По своей логике данная модель близка к режиму пользовательской загрузки фонограмм, реализованному в отдельных цифровых музыкальных сервисах. Однако с практической точки зрения она существенно усложняет клиентский опыт: пользователю необходимо самостоятельно находить, загружать и поддерживать музыкальную библиотеку. Кроме того, сохраняется риск негативной реакции со стороны бизнеса в случае возникновения претензий, связанных с отсутствием необходимых договоров с РАО и ВОИС, поскольку требования к лицензированию публичного воспроизведения музыки не являются очевидными для большинства непрофессиональных пользователей.

Третий вариант представляет собой модификацию предыдущей модели, при которой функции по наполнению музыкальной библиотеки передаются внешнему поставщику или иному третьему лицу. Формально сам сервис в этом случае также может оставаться инструментом рекомендаций и воспроизведения, не заявляя себя в качестве правообладателя или дистрибьютора музыкального каталога. Однако подобная конструкция оказывается юридически менее прозрачной, поскольку вопрос о легальности происхождения файлов и объеме передаваемых прав фактически выносится за пределы основной платформы. В результате такая схема находится в правовой «серой зоне».

Наконец, четвертый вариант предполагает развитие сервиса по модели white label или B2B-модуля, интегрируемого в уже существующие решения для музыкального сопровождения бизнеса. В такой архитектуре проектируемый продукт не берет на себя функцию предоставления музыкального каталога, а отвечает за интеллектуальный слой сервиса: сбор сигналов о предпочтениях посетителей, формирование групповых рекомендаций и передачу рекомендованного контента во внешнюю систему воспроизведения. Данный подход представляется наиболее институционально устойчивым, поскольку он позволяет встроить технологию коллективных рекомендаций в уже существующую инфраструктуру легальной музыки для бизнеса. Показательно, что «Яндекс» уже использует партнерскую модель в сегменте B2B-музыки, сотрудничая с Cubic Music в рамках продукта «Яндекс Станция. Музыка для бизнеса».

С учетом обозначенных ограничений целесообразно рассматривать наличие действующих договоров с РАО и ВОИС как обязательное условие подключения функций групповых музыкальных рекомендаций для бизнеса. Более того, в рамках полноценного B2B-сервиса заключение необходимых лицензионных соглашений, сопровождение взаимодействия с организациями по коллективному управлению правами и подготовка отчетных документов могут быть включены в минимальный пакет услуг платформы. Такой подход снижает правовые риски для клиентов, делает модель использования сервиса более прозрачной и одновременно усиливает его ценностное предложение для бизнеса, поскольку снимает с последнего значительную часть административной и юридической нагрузки.

\section{Исследование потребностей, проблем и ожиданий целевых сегментов}
\label{sec:market:needs}

Для уточнения потребностей целевых сегментов в рамках настоящей работы было проведено качественное исследование в формате экспертных глубинных интервью. В интервью участвовали представители индустрий, выделенных в разделе~\ref{sec:market:segmentation}: представители условно ресторанного и барного сегмента, организаторы мероприятий, DJ, а также юристы, работающие с музыкальными правами. Целью этого этапа было не получение статистически репрезентативной картины, а проверка продуктовых гипотез и выявление практических болей, которые не всегда видны из анализа рынка и публичных описаний сервисов.

Полученные интервью интерпретируются как качественная валидация проблемного поля. Поэтому в данном подразделе не делаются выводы о частоте тех или иных практик на всем рынке. Вместо этого фиксируются устойчивые мотивы, которые повторялись в разговорах с экспертами и хорошо согласуются с выводами предыдущих разделов: музыка воспринимается бизнесом одновременно как элемент атмосферы, операционная задача и источник правового риска; для мероприятий музыка является частью сценария, а не просто фоном; для DJ важны реакция аудитории и гибкость подготовки сета; для посетителей потенциальная ценность персонализации должна быть уравновешена прозрачной моделью согласия и приватности.

\subsection{Заведения}
\label{subsec:market:needs-business}

Представители ресторанного и барного сегмента в первую очередь выделяли не алгоритмическую, а правовую проблему музыкального сопровождения. Для кафе, баров, ресторанов и клубов одним из наиболее заметных страхов является риск привлечения к ответственности за публичное воспроизведение музыки без необходимых разрешений. Этот страх усиливается за счет публичных кейсов о взыскании компенсаций и новостей о росте претензионной активности правообладателей. Так, по данным «Коммерсанта», Российское авторское общество за 2025 год подало 2,3 тыс. исков на 515,5 млн руб.; сумма требований выросла на 42,6\% по сравнению с предыдущим годом~\cite{kommersant_rao_claims_2026}. Этот источник не описывает конкретно целевую аудиторию предлагаемого сервиса, но подтверждает общий фон правовой напряженности вокруг публичного воспроизведения музыки.

Интервью показали, что проблема не сводится к нежеланию бизнеса платить за музыку. Часть представителей малого и среднего бизнеса просто недостаточно понимает, как легально оформить публичное воспроизведение, какие организации участвуют в процессе и чем отличаются права на произведение от прав на фонограмму. Как было показано в разделе~\ref{sec:market:legal}, публичное использование записанной музыки может затрагивать как авторские, так и смежные права, а отчетность перед РАО и ВОИС требует фиксации фактически использованного репертуара~\cite{rao_report_forms,vois_report_form}. Для непрофильного бизнеса эта схема выглядит сложной и плохо операционализируемой.

Отдельный мотив интервью связан с перекладыванием ответственности между участниками события. В небольших мероприятиях может быть неочевидно, кто именно должен обеспечить правомерность воспроизведения: площадка, организатор, приглашенный DJ или подрядчик, отвечающий за звук. Опрошенные юристы указывали, что претензии часто адресуются площадке, поскольку именно она является наиболее очевидным и доказуемым ответчиком: нарушение происходит на ее территории, а сам факт воспроизведения легче зафиксировать в пространстве заведения. Со стороны артистов и DJ, напротив, была описана лишь единичная практика предоставления списка исполненных треков на крупном мероприятии, причем не для предварительного согласования с организатором, а как элемент последующей отчетности перед правообладателями.

Из этого следует важный продуктовый вывод: сервис коллективных музыкальных рекомендаций не должен восприниматься бизнесом как еще один интерфейс выбора треков. Для целевого сегмента существенной частью ценности становится снижение юридической и административной неопределенности. Встроенная отчетность, сопровождение взаимодействия с РАО и ВОИС или интеграция с уже легализованным B2B-контуром могут быть не дополнительной функцией, а условием доверия к продукту. Иначе бизнес получает новый рекомендательный слой, но сохраняет прежний страх перед правовыми последствиями.

В интервью также проявилась роль посредников. Некоторые заведения передают музыкальное оформление сторонним организациям, которые уже от имени клиента взаимодействуют с поставщиками музыки для бизнеса и помогают поддерживать нужный стиль звучания. Это означает, что среди потенциальных целевых аудиторий следует учитывать не только конечные площадки, но и B2B-подрядчиков, оказывающих услуги музыкального и маркетингового оформления. Для них сервис может стать не заменой текущей деятельности, а технологическим модулем, повышающим качество подбора музыки и упрощающим отчетность.

Еще один важный вывод связан с различием режимов использования музыки внутри одного и того же заведения. Управляющий бара описывал, что в периоды низкого потока клиентов музыка выполняет функцию заполнения тишины и поддержания общего фона. Однако в пиковые вечера, особенно по пятницам и субботам, музыка становится частью событийного опыта: заведение нанимает DJ или использует популярные плейлисты, потому что гости приходят не только пообщаться, но и танцевать. Следовательно, один и тот же бизнес может нуждаться в разных музыкальных режимах: фоновом, брендированном, танцевальном, событийном. Это напрямую поддерживает продуктовую логику главы~\ref{ch:service-concept}, где сервис должен позволять бизнесу задавать рамки по стилю, энергичности и сценарию использования.

С учетом этих наблюдений ранняя целевая аудитория сервиса, вероятно, не совпадает с бизнесами, которые ищут самый дешевый способ «закрыть тишину» и минимально легализовать публичное воспроизведение. Более перспективными выглядят площадки, для которых критично использование популярной или узнаваемой музыки и которые готовы поддерживать более сложный правовой и продуктовый контур: клубы, бары с танцполом, концертные и event-площадки, а также заведения, где музыка является заметной частью позиционирования. Для таких клиентов ценность сервиса состоит не только в фоновой легальности, но и в более точном попадании в аудиторию, находящуюся в пространстве здесь и сейчас.

\subsection{Организаторы мероприятий и диджеи}
\label{subsec:market:needs-events}

Интервью с организаторами мероприятий и DJ показали, что событийный сегмент существенно отличается от типичного рынка музыки для бизнеса. Для вечеринок и клубных событий часто недостаточно каталога фоновой легальной музыки: организаторам нужны узнаваемые треки, понятные аудитории, а также возможность работать с эстетикой конкретного жанра, сцены или комьюнити. В таких сценариях музыкальный выбор обычно остается за DJ, тогда как организатор или арт-директор задает рамки: жанр, общее настроение, уровень энергии, формат вечера и референсные сеты.

Референсные плейлисты и DJ-сеты выполняют в этом контексте особую функцию. Они не обязательно являются списком треков, которые нужно воспроизвести буквально. Напротив, часто это способ объяснить направление: «нужно в похожем стиле, но не эти же композиции». Следовательно, для сервиса коллективных рекомендаций важна возможность переводить такие референсы в машинно читаемый вид: через сопоставление треков с каталогом платформы, анализ аудиоэмбеддингов или поиск близких композиций в пространстве музыкальных представлений. Эта логика согласуется с технической частью работы, где аудиоэмбеддинги YAMBDA используются как содержательный сигнал для группового агрегатора.

Отдельное место в интервью заняла структура вечеринки. Типичный лайнап строится вокруг кульминации: ближе к середине или началу второй половины события выступает хедлайнер, ради которого приходит значительная часть аудитории. До него играют артисты, задача которых состоит в постепенном «разогреве» зала, а после него музыкальная интенсивность может поддерживаться или снижаться в зависимости от сценария. Для opening DJ важна не максимальная энергичность, а аккуратная подготовка аудитории: нужно создать движение, не перегреть зал слишком рано и не разрушить ожидание основного выступления. Ограничения могут задаваться через BPM, жанр, уровень энергии или более неформальное описание концепции вечеринки.

В этом контексте сервис коллективных рекомендаций может быть рассмотрен как инструмент поддержки opening-сценария. Он потенциально способен учитывать предпочтения уже присутствующих гостей и выбирать треки, которые мягко повышают вовлеченность, оставаясь внутри рамок, заданных организатором. При этом важно не смешивать такую поддержку с полной заменой DJ. Живой исполнитель учитывает не только музыкальные предпочтения, но и динамику зала, сценическое поведение, драматургию переходов и невербальную реакцию аудитории. Полевые исследования электронной танцевальной сцены также показывают, что поведение DJ и его взаимодействие с залом связаны с положительным аффектом аудитории~\cite{rudnicki2025djs}. Следовательно, алгоритм может быть полезным источником сигнала, но не исчерпывает профессиональную роль исполнителя.

Содержательно наиболее интересным оказался инсайт о так называемой «теплой» аудитории. DJ описывали ситуацию, когда часть гостей уже готова переместиться к танцполу, но еще остается у бара, за столами или в промежуточных зонах пространства. В такие моменты задача состоит не просто в повышении BPM, а в поиске музыкальных «крюков»: знакомых фрагментов, треков или стилистических решений, которые способны мягко вовлечь именно эту группу людей. Знание музыкальных предпочтений присутствующих пользователей теоретически дает для такой задачи дополнительный инструмент: система может подсказать, какие треки с большей вероятностью окажутся знакомыми или приемлемыми для людей, находящихся в менее вовлеченной зоне.

Этот вывод напрямую связан с перспективами, описанными в главе~\ref{ch:prototype}. Если в пространстве используется несколько BLE-маяков, то система может приблизительно различать посетителей у танцпола, бара или в лобби; при этом, как отмечено в подразделе~\ref{subsec:ch6:multi-ble}, такой подход нельзя трактовать как точную геолокацию. Для продуктовой гипотезы достаточно более слабого сигнала: пользователь находится ближе к одной зоне, чем к другой. Тогда коллективная рекомендательная модель может не только формировать общий музыкальный поток, но и помогать работать с разной степенью вовлеченности аудитории в пространстве.

Интервью с DJ также позволили уточнить, почему сценарий ассистента для DJ требует отдельной архитектуры. В подготовке сета можно выделить несколько практических подходов. Первый предполагает заранее собранный сет, подготовленные переходы и контрольные точки в специализированном программном обеспечении. Такой подход дает высокий контроль над качеством сведения, но снижает гибкость реакции на аудиторию. Второй подход состоит в подготовке большой библиотеки треков и переходов, из которой DJ выбирает материал в ходе выступления. Третий подход ближе к импровизации: исполнитель работает с максимально широкой медиатекой и принимает решения в моменте, что требует большого опыта, быстрой навигации по библиотеке и уверенного технического навыка. Современные DJ-платформы действительно позиционируются как инструменты для подготовки и управления музыкальными библиотеками и сетами~\cite{rekordbox_software}.

Общее между этими подходами состоит в том, что профессиональная музыкальная библиотека DJ часто существует как набор локальных файлов, а не как плейлист внутри массового стримингового сервиса. Это ограничивает прямое применение рекомендательной модели, работающей с каталогом «Яндекс Музыки». Для сценария DJ-ассистента возможны два технических пути. Первый~--- сопоставление локальных файлов с каталогом платформы по названию, метаданным и цифровым аудиоотпечаткам; если трек найден, система может оценивать его релевантность текущей аудитории через уже имеющиеся каталожные идентификаторы. Второй~--- работа через аудиохарактеристики и аудиоэмбеддинги, когда точное сопоставление с каталогом невозможно.

Второй путь особенно важен из-за практик, описанных респондентами. DJ могут специально искать музыку, которая плохо распознается массовыми сервисами, использовать редкие загрузки с открытых площадок, неопубликованные версии, edit-версии, ремиксы и мэшапы. Такие треки повышают уникальность сета, но затрудняют идентификацию по названию или аудиоотпечатку. В этом случае аудиоэмбеддинги могут помочь найти содержательно близкий музыкальный материал или оценить близость локального трека к групповому профилю, но такой подход будет менее надежным, чем работа с точными каталожными идентификаторами и коллаборативной историей прослушиваний.

Следовательно, ассистент для DJ должен рассматриваться не как простое расширение бизнес-плеера, а как отдельный продуктовый сценарий. Он должен уметь работать с локальными файлами, поддерживать сопоставление с каталогом, использовать аудиопризнаки, учитывать ограничения по стилю и не мешать профессиональной драматургии сета. В текущей работе этот сценарий не реализуется как отдельный прототип, но он задает важное направление развития технологии, подробнее раскрытое в подразделе~\ref{subsec:ch6:dj-assistant}.

\subsection{Ключевые выводы}
\label{subsec:market:needs-users}
\label{subsec:market:needs-insights}

Хотя проведенные экспертные интервью были сфокусированы преимущественно на профессиональных участниках индустрии, продуктовая концепция сервиса невозможна без учета позиции посетителя. С точки зрения пользователя ценность состоит в возможности мягко влиять на музыкальную среду без явного заказа треков, публичного голосования или конфликта с другими гостями. Такая логика отличается от visitor ordering: пользователь не покупает право поставить конкретную композицию, а его музыкальный профиль становится одним из вкладов в общий групповой поток.

Одновременно именно пользовательский сценарий несет наиболее чувствительный приватностный риск. История прослушиваний может восприниматься как личная информация, а факт нахождения в заведении или на мероприятии~--- как контекст, который не должен автоматически становиться доступным бизнесу. Поэтому сервис должен быть построен вокруг явного информированного согласия, понятного объяснения пользы и возможности отказаться от участия. Как сформулировано в главе~\ref{ch:service-concept}, владелец заведения должен получать только музыкальный результат и агрегированные показатели, но не персональные данные посетителей.

Таким образом, ожидания целевых сегментов оказываются разнонаправленными. Бизнесу нужна управляемость, правовая определенность и минимум операционной сложности. Организатору мероприятия и DJ нужен дополнительный сигнал об аудитории, но без потери художественного и сценарного контроля. Пользователю нужна понятная ценность и защита приватности. Платформе необходимо обеспечить такую архитектуру, в которой рекомендательный слой, лицензирование, отчетность и пользовательское согласие не противоречат друг другу.

По итогам экспертных интервью можно сформулировать несколько ключевых выводов. Для заведений музыкальное сопровождение является не только вопросом атмосферы, но и источником правового риска. Страх претензий РАО, ВОИС и правообладателей выступает самостоятельным фактором принятия решений, а значительная часть боли бизнеса связана с непониманием процедуры легализации публичного воспроизведения и распределения ответственности между площадкой, организатором и исполнителем. Поэтому автоматическая отчетность и сопровождение лицензионного контура могут быть не дополнительной функцией, а базовой частью ценностного предложения сервиса.

Наиболее перспективными ранними клиентами выглядят площадки, для которых важна популярная и узнаваемая музыка: клубы, бары с танцполом, концертные пространства и организаторы мероприятий. Для событийного сегмента критична не просто генерация плейлиста, а поддержка музыкальной драматургии: разогрев аудитории, работа с уровнем энергии и соблюдение жанровых рамок. В этом контексте сервис может быть полезен DJ как ассистент и источник дополнительного сигнала о предпочтениях аудитории, но не должен позиционироваться как замена профессионального исполнителя.

Отдельный вывод связан с техническими ограничениями DJ-сценария. Профессиональные DJ-библиотеки часто состоят из локальных файлов, редких версий, ремиксов и мэшапов, поэтому сценарий DJ-ассистента требует механизмов сопоставления треков и анализа аудиоэмбеддингов. Пользовательская ценность сервиса, в свою очередь, возникает только при прозрачном добровольном участии: посетитель должен понимать, как он влияет на музыку, какие данные используются и почему площадка не получает его персональный профиль.

\section{Основные барьеры и ограничения внедрения сервиса коллективных музыкальных рекомендаций}
\label{sec:market:barriers}

Первый и наиболее чувствительный барьер связан с приватностью. Сервис предполагает определение факта присутствия пользователя в заведении или на мероприятии и использование истории его музыкальных предпочтений для формирования группового плейлиста. Даже если владелец заведения не получает персональные данные, сам факт участия пользователя в таком сценарии требует информированного согласия и прозрачной коммуникации. Пользователь должен понимать, что он влияет на музыкальный поток, какие данные используются и как отказаться от участия.

Второй барьер связан с правовой и лицензионной архитектурой. Как показано в разделе~\ref{sec:market:legal}, договоры с РАО и ВОИС решают только часть проблемы публичного воспроизведения. Для промышленной реализации сервиса на базе массового каталога требуется отдельное B2B-урегулирование прав на использование фонограмм через платформу. Без такого контура сервис может остаться лишь рекомендательной надстройкой, которая формирует плейлист, но не обеспечивает полностью легальную цепочку воспроизведения.

Третий барьер имеет организационный характер. Для малого и среднего бизнеса музыка часто является второстепенной операционной задачей, а не самостоятельным объектом управления. Даже если сервис способен улучшить музыкальную среду, заведению нужно установить оборудование или настроить устройство воспроизведения, обучить персонал, назначить ответственного за профиль, задать ограничения и регулярно проверять корректность работы. Следовательно, успешный продукт должен минимизировать операционные действия клиента и предлагать простую модель запуска.

Четвертый барьер связан с технологией фиксации присутствия. BLE-beacon позволяет автоматизировать определение нахождения пользователя в зоне заведения, но требует совместимости смартфона, разрешений приложения, стабильной работы маяков и аккуратной настройки зон. Альтернативой или дополнением может быть QR-код, который проще объяснить пользователю и использовать в промо-сценариях, но хуже подходит для пассивного и динамического обновления состава аудитории. Поэтому в концепции сервиса целесообразно рассматривать BLE как основной автоматизированный механизм, а QR~--- как резервный или стартовый сценарий.

Пятый барьер~--- зависимость от экосистемы Яндекса. Сильная сторона проектируемого сервиса заключается в том, что «Яндекс Музыка» уже обладает пользовательскими профилями и историей прослушивания. Но это же является ограничением: посетители, не использующие сервис, не смогут в полной мере участвовать в формировании потока, а качество рекомендаций будет зависеть от полноты и актуальности их музыкальной истории. Для бизнеса это означает, что продуктовая ценность будет выше в аудиториях с высокой проникновенностью «Яндекс Музыки».

Наконец, существуют репутационные и пользовательские риски. Бизнес может опасаться потери контроля над атмосферой, если музыка начнет «подстраиваться под всех» без ограничений. Пользователь может опасаться слежения или нежелательного раскрытия вкусов. Эти риски снижаются архитектурой, в которой заведение задает рамки, стоп-листы и уровень энергичности, а пользовательские данные остаются внутри платформы. Однако полностью устранить их нельзя; поэтому коммуникация приватности и контрольных механизмов является обязательной частью сервиса.

\section{Оценка рыночного потенциала сервиса}
\label{sec:market:potential}

Оценка рыночного потенциала проектируемого сервиса должна строиться уже не от всей совокупности заведений, где в принципе может звучать музыка, а от тех форматов, где музыкальная среда является значимой частью клиентского опыта. Как показано в предыдущих разделах, сервис коллективных музыкальных рекомендаций не сводится к легальному фоновому плееру: его ценность возникает в ситуации, когда плейлист может быть адаптирован под фактический состав присутствующей аудитории. Поэтому при оценке рынка необходимо исключить значительную часть организаций, для которых музыка является лишь нейтральным фоном и не образует самостоятельной продуктовой ценности.

В качестве релевантного ядра рассматриваются прежде всего музыкально чувствительные форматы: вечерние бары, бары с танцами и DJ-сетами, ночные клубы, заведения с танцполом, часть restaurant-bar и lounge-форматов, а также отдельные событийные сценарии. Обычные кафе, кофейни, магазины, фитнес-клубы, гостиницы и салоны красоты могут регулярно использовать музыку, однако в большинстве таких случаев эффект от динамической групповой персонализации слабее, чем в заведениях, где гости приходят не только за базовой услугой, но и за атмосферой. Такой подход согласуется с выводами разделов~\ref{sec:market:music-service}--\ref{sec:market:barriers}: важен не сам факт звучания музыки, а ее соответствие аудитории, сценарию посещения и роли заведения.

Методика оценки строится через сочетание двух контуров. Первый контур исходит из уже существующего легального рынка пользователей РАО и ВОИС. Он отражает наиболее зрелую часть бизнеса: организации, которые признают необходимость легальных выплат за публичное использование музыки и потому потенциально ближе к покупке дополнительного технологического слоя. Второй контур строится от числа потенциально релевантных заведений в интересующих сегментах и проверяет, какая часть таких заведений по масштабу выручки в принципе могла бы нести расходы на легальное музыкальное сопровождение. Эти подходы нельзя механически складывать: часть заведений одновременно входит и в легальный контур, и в расчетную базу релевантных индустрий. Поэтому итоговая оценка TAM, SAM и SOM ниже задается как сценарная триангуляция, а не как точная статистическая перепись рынка.

Для экономической фильтрации рынка важно сначала определить типичный масштаб платежей за легальное публичное воспроизведение музыки. В подготовленном исследовании для этой цели была собрана малая публичная выборка площадок с танцполом в восьми крупных городах России. В нее вошли опубликованные карточки и страницы заведений; для объектов с указанием «более 500 м\textsuperscript{2}» использовалась нижняя граница, а для аномально крупных площадок применялась винзоризация на уровне 800 м\textsuperscript{2}, чтобы единичные клубно-ресторанные комплексы не смещали оценку вверх. Итоговая репрезентативная площадь составила около 380 м\textsuperscript{2}. Это не официальная средняя по отрасли, а расчетный ориентир для типового городского заведения с танцполом, пригодный для рыночного моделирования. Городские оценки, на которых основан этот ориентир, приведены на Рисунке~\ref{fig:market:venue-area}.

\begin{figure}[htbp]
\centering
\begin{tikzpicture}
\begin{axis}[
  ybar,
  width=0.95\textwidth,
  height=0.42\textwidth,
  ymin=0,
  ymax=700,
  ylabel={м\textsuperscript{2}},
  symbolic x coords={Мск,СПб,Нск,Екб,Кзн,НН,Сам,Рнд},
  xtick=data,
  nodes near coords,
  nodes near coords style={font=\scriptsize},
  tick label style={font=\small},
  ymajorgrids=true,
  grid style={gray!25},
  bar width=15pt,
  enlarge x limits=0.08
]
\addplot+[fill=gray!35,draw=gray!70] coordinates {
  (Мск,370)
  (СПб,350)
  (Нск,385)
  (Екб,380)
  (Кзн,545)
  (НН,180)
  (Сам,665)
  (Рнд,300)
};
\end{axis}
\end{tikzpicture}
\caption{Оценка средней площади заведений с танцполом по городам}
\label{fig:market:venue-area}
\end{figure}

Для дальнейшего расчета эта площадь используется как базовый сценарий. При этом сама правовая природа публичного воспроизведения музыки подробно рассмотрена в разделе~\ref{sec:market:legal}; здесь платежи РАО и ВОИС используются не для повторения правового анализа, а как экономический индикатор того, какую регулярную нагрузку несет бизнес, работающий с узнаваемой музыкой в легальном контуре. По действующей тарифной логике для предприятия общественного питания с озвучиваемой площадью 380 м\textsuperscript{2} платеж рассчитывается по ставке 38~руб. за 1 м\textsuperscript{2}: 14\,440~руб. в месяц по линии РАО и 14\,440~руб. по линии ВОИС, всего 28\,880~руб. в месяц. Если же объект по своему позиционированию и фактическому режиму работы подпадает под клубный тариф, то при площади 380 м\textsuperscript{2} применяется фиксированная ставка 17\,000~руб. в месяц у каждой из организаций, всего 34\,000~руб. в месяц~\cite{rao_tariffs_2024,vois_tariffs_2024}. Сравнение двух базовых сценариев показано на Рисунке~\ref{fig:market:rao-vois-payment}.

\begin{figure}[htbp]
\centering
\begin{tikzpicture}
\begin{axis}[
  ybar,
  width=0.72\textwidth,
  height=0.42\textwidth,
  ymin=0,
  ymax=36000,
  ylabel={руб./мес.},
  symbolic x coords={Общепит,Клуб},
  xtick=data,
  nodes near coords,
  nodes near coords style={font=\small,/pgf/number format/1000 sep={\,}},
  tick label style={font=\small},
  ymajorgrids=true,
  grid style={gray!25},
  bar width=34pt,
  enlarge x limits=0.35,
  scaled y ticks=false,
  y tick label style={/pgf/number format/1000 sep={\,}}
]
\addplot+[fill=gray!35,draw=gray!70] coordinates {
  (Общепит,28880)
  (Клуб,34000)
};
\end{axis}
\end{tikzpicture}
\caption{Суммарный ежемесячный платеж РАО и ВОИС при площади 380 м\textsuperscript{2}}
\label{fig:market:rao-vois-payment}
\end{figure}

Разница между этими сценариями показывает, что для заведений с танцполом ключевым фактором является не только площадь, но и тарифная квалификация объекта. Для меньших объектов платежи пропорционально ниже: например, бар площадью 120 м\textsuperscript{2} при ставке 40~руб. за 1 м\textsuperscript{2} у каждой организации дает совокупную нагрузку около 9\,600~руб. в месяц. Для более крупного restaurant-bar площадью 500 м\textsuperscript{2} при ставке 38~руб. за 1 м\textsuperscript{2} совокупный платеж составит около 38\,000~руб. в месяц. Клубные тарифы меняются ступенчато, поэтому при росте площади нагрузка увеличивается не линейно, а при переходе между тарифными диапазонами.

В настоящей работе принимается расчетное допущение: расходы на платежи РАО и ВОИС должны составлять не более 0{,}5\% ежемесячной выручки заведения. Этот порог не является нормативным стандартом и не претендует на универсальное описание экономики заведений. Он используется как консервативный фильтр, позволяющий отделить бизнесы, для которых регулярные легальные платежи за музыку выглядят экономически переносимыми, от объектов, где такая нагрузка, вероятно, воспринималась бы как чрезмерная. Минимальная месячная выручка может быть описана формулой

\begin{equation*}
    R_{\min} = \frac{P_{\mathrm{RAO+VOIS}}}{0{,}005},
\end{equation*}

где $P_{\mathrm{RAO+VOIS}}$~--- совокупный ежемесячный платеж РАО и ВОИС, а $R_{\min}$~--- минимальная расчетная выручка заведения. Для базового сценария общепита при платеже 28\,880~руб. необходимая выручка составляет 5\,776\,000~руб. в месяц. Для клубного сценария при платеже 34\,000~руб. минимальная выручка равна 6\,800\,000~руб. в месяц. Для бара площадью 120 м\textsuperscript{2} соответствующий порог составит около 1\,920\,000~руб., а для объекта площадью 500 м\textsuperscript{2}~--- около 7\,600\,000~руб. Эти расчеты не задают цену будущей подписки, а только определяют экономический фильтр для оценки платежеспособной части рынка.

Первый расчетный контур связан с уже легализованным использованием музыки. По подготовленным материалам на конец 2024 года у РАО было 26\,130 договоров с пользователями, а у ВОИС~--- 21\,684 действующих договора; существенная часть сборов ВОИС относилась именно к публичному исполнению фонограмм~\cite{rao_annual_report_2024,vois_annual_report_2024}. Эти показатели нельзя напрямую интерпретировать как число физических заведений: один договор может покрывать несколько точек сетевого бизнеса, а часть договоров относится к форматам, не являющимся целевыми для рассматриваемого сервиса. Поэтому данные отчетов используются как нижний ориентир рыночного ядра, а не как готовая оценка TAM.

Для перехода от договоров к релевантным точкам применяются несколько сужающих допущений. Во-первых, из общего массива пользователей выделяются только те, кто связан с публичным воспроизведением в коммерческих и досуговых пространствах. Во-вторых, из этого массива отбираются не все пользователи музыки, а только музыкально чувствительные форматы: бары, клубы, restaurant-bar, lounge и площадки с выраженной вечерней программой. В-третьих, на результат накладывается фильтр 0{,}5\% выручки, поскольку не каждый пользователь легального контура имеет достаточный масштаб для покупки дополнительного сервиса поверх обязательных правообладательских платежей. В таком прочтении уже существующий легальный контур дает консервативное ядро порядка 3--4 тыс. потенциально релевантных точек, а при более широкой трактовке сетевых объектов и nightlife-форматов~--- до 5 тыс. точек.

Второй расчетный контур строится «снизу вверх» от числа объектов в интересующих индустриях. В качестве открытых ориентиров используются оценки Контур.Фокуса, согласно которым в России насчитывалось 10\,902 бара и 158\,902 ресторана, а также отраслевые оценки более чем 10 тыс. кальянных клубов и сигарных lounge-форматов~\cite{kontur_foodservice_2025,retail_hookah_lounge_2026}. Для ночных клубов открытая федеральная статистика менее надежна, поэтому их количество задается сценарно в диапазоне 1--2 тыс. объектов. Однако все эти базы не могут целиком быть отнесены к целевому рынку. В TAM включается только та доля форматов, где музыка заметна, узнаваема и потенциально требует адаптации под текущую аудиторию: в базовом сценарии около 65\% баров, 4\% ресторанов, 25\% lounge-форматов и практически все клубы. Такой расчет дает около 17--18 тыс. точек в базовом сценарии; при более жестких или более широких коэффициентах разумный коридор TAM составляет 14--23 тыс. точек.

Следующий шаг состоит в сужении TAM до SAM. Здесь применяется одновременно экономический фильтр 0{,}5\% и правовая адресуемость рынка. Порог 0{,}5\% делает оценку существенно строже, чем простое включение всех вечерних заведений: небольшие бары и lounge-объекты могут быть музыкально релевантными, но не иметь выручки, при которой совокупные выплаты РАО и ВОИС выглядят комфортной регулярной нагрузкой. Поэтому в базовом сценарии в SAM включается не весь TAM, а только часть заведений, которые либо уже находятся в легальном контуре, либо по масштабу бизнеса потенциально способны перейти в него. Для баров это примерно треть релевантного TAM-среза с дополнительной поправкой на правовую готовность; для restaurant-bar~--- еще более узкая доля, поскольку большая часть ресторанного рынка не относится к музыкально чувствительным форматам; для lounge-форматов~--- ограниченный сегмент наиболее вечерних и атмосферных заведений; для клубов~--- более высокая доля, так как музыка составляет ядро продукта. В совокупности такой пересчет дает базовую оценку около 3{,}5--4 тыс. точек и сценарный диапазон SAM примерно 3--5 тыс. точек.

Событийный рынок в этой модели рассматривается отдельно. Мероприятие отличается от постоянной точки единицей учета, циклом продажи, бюджетодержателем и режимом лицензирования: для события важен бюджет конкретного вечера или билетная выручка, а не месячная выручка площадки. Поэтому события не включаются механически в точечные оценки TAM, SAM и SOM. Они могут стать отдельным каналом применения сервиса, особенно для клубных вечеринок, корпоративных событий и брендированных мероприятий, но для целей финансового моделирования их следует считать отдельным сценарием, чтобы не смешивать подписочную модель постоянных заведений и проектную модель разовых событий.

Итоговая оценка постоянного рынка может быть сформулирована следующим образом. TAM проектируемого сервиса составляет около 14--23 тыс. музыкально чувствительных точек, базово~--- 17--18 тыс. точек. SAM, то есть экономически и правово адресуемая часть этого рынка, составляет около 3--5 тыс. точек, базово~--- порядка 4 тыс. точек. SOM в рамках настоящей работы целесообразно трактовать не как размер первого технического пилота, а как достижимый объем раннего коммерческого масштабирования после проверки продукта на ограниченном числе площадок. При такой трактовке учитывается, что сервис потенциально может распространяться не только через прямые продажи отдельным независимым заведениям, но и через партнерства с B2B-поставщиками музыки, интеграции с существующими решениями для заведений, экосистемные каналы Яндекса и сетевых клиентов, где один договор может охватывать несколько физических точек. С учетом этих предпосылок реалистично достижимой долей SAM можно считать около 10--14\% базового адресуемого рынка, что соответствует диапазону 450--550 заведений. Базовый ориентир SOM в дальнейшем финансовом моделировании может быть принят на уровне около 500 заведений. Такая оценка остается сценарной и не означает гарантированного охвата рынка, однако лучше отражает коммерческий горизонт продукта после успешного подтверждения ценности сервиса.

Таким образом, рынок сервиса коллективных музыкальных рекомендаций существует не в масштабе «всех заведений с музыкой», а в масштабе конкретного сегмента музыкально чувствительных площадок, способных нести расходы на легальное музыкальное сопровождение. Полученные оценки показывают, что потенциальный рынок достаточно узок, чтобы требовать точного позиционирования и аккуратного выбора первых клиентов, но достаточно велик для пилотирования и последующей коммерциализации. На следующем этапе эти значения могут использоваться как база для расчета финансовой модели и возможной стоимости ежемесячной подписки, необходимой для окупаемости проекта, однако сама цена подписки в рамках данного подраздела не задается.

\section{Выводы}
\label{sec:market:conclusion}

Проведенный анализ показывает, что музыкальное сопровождение заведений и мероприятий является самостоятельной сферой услуг, в которой музыка выполняет не только эстетическую, но и сервисную, маркетинговую и сценарную функцию. В литературе наиболее устойчивые положительные эффекты связываются не с простым фактом наличия музыки, а с ее субъективной привлекательностью и уместностью для контекста. Это создает теоретическое основание для гипотезы о ценности адаптивного музыкального потока, учитывающего предпочтения реально присутствующей аудитории.

Российский рынок музыкального сопровождения неоднороден. В нем сосуществуют локальные плейлисты и потребительские стриминги, брендированные музыкальные решения, специализированные B2B-сервисы и профессиональные событийные практики, связанные с работой DJ и музыкальных редакторов. Конкурентный анализ показывает, что существующие сервисы в основном закрывают легальность, каталоги, управление воспроизведением, аудиомаркетинг и иногда заказ треков гостями. При этом почти отсутствуют решения, которые используют предпочтения фактически присутствующих посетителей для мягкого коллективного формирования музыкального потока внутри рамок заведения.

Правовая среда формирует для такого сервиса существенные ограничения. Для публичного воспроизведения записанной музыки бизнесу, как правило, необходимо учитывать как авторские права на произведение, так и смежные права на фонограмму, что предполагает взаимодействие с РАО и ВОИС. Кроме того, потребительская подписка на музыкальный стриминг не равна B2B-праву на публичное коммерческое использование. Поэтому промышленная версия сервиса должна включать не только рекомендательную модель, но и устойчивый лицензионный и отчетный контур.

Рыночная оценка показывает, что наиболее релевантным адресуемым рынком для старта является не весь массив заведений, где в принципе может звучать музыка, а сегмент музыкально чувствительных форматов: бары, клубы, restaurant-bar, lounge-объекты и площадки с выраженной вечерней программой. В сценарной оценке TAM такого рынка составляет около 14--23 тыс. точек, SAM~--- около 3--5 тыс. точек, а реалистично достижимый SOM на горизонте раннего коммерческого масштабирования~--- 450--550 точек. Событийный сценарий следует рассматривать отдельно, поскольку мероприятия отличаются от постоянных заведений единицей учета, циклом продажи и логикой лицензирования.

С учетом выявленных предпосылок разработка сервиса коллективных музыкальных рекомендаций представляется перспективной не как замена существующих B2B-плееров, а как новый слой пользовательского и продуктового опыта: легальный, управляемый, приватный и способный учитывать реальный состав аудитории. Концепция такого сервиса раскрывается в следующей главе.

% TODO(групповая ВКР): указать распределение задач между участниками проекта в рамках главы 1.
